% SOURCE_STATEMENT: "For general null cones, characterize convergence of the Hawking energy, compare its limit with the Trautman–Bondi mass, and derive the corresponding null Penrose inequality."
% Preparation and final self-audit: CONTRIBUTING.md prompts C and D; source/audit artifacts retained in the local physics_statements_2026_10_11 batch.
% !TeX program = lualatex
\documentclass[11pt]{article}
\usepackage[margin=25mm]{geometry}
\usepackage{fontspec}
\setmainfont{Latin Modern Roman}
\usepackage{amsmath,amssymb,mathtools}
\usepackage{unicode-math}
\setmathfont{Latin Modern Math}
\usepackage{xurl,hyperref}
\hypersetup{colorlinks=true,urlcolor=blue}
\setcounter{secnumdepth}{0}
\setlength{\parindent}{0pt}
\setlength{\parskip}{5pt}
\setlength{\emergencystretch}{3em}
\newcommand{\sref}[2]{\hyperlink{source:#1:#2}{[S#2]}}
\newcommand{\cataloguescope}[1]{\paragraph{Recorded target.}#1}
\newcommand{\R}{\mathbb R}
\newcommand{\C}{\mathbb C}
\newcommand{\N}{\mathbb N}
\newcommand{\Z}{\mathbb Z}
\newcommand{\Q}{\mathbb Q}
\newcommand{\F}{\mathbb F}
\newcommand{\eps}{\varepsilon}
\begin{document}
\section{Hawking Energy Limits and the Null Penrose Inequality}\label{statement}

\subsection{Definitions and mathematical statement}

\paragraph{Shared notation.}$\mathbb N=\{1,2,\ldots\}$; the standard number fields have their usual meanings. All additional conventions are stated locally.


\paragraph{Definitions and conventions.}
A null cone is a hypersurface generated by null geodesics; its spacelike two-surface sections carry null expansions. Hawking energy is a quasilocal mass constructed from a section\textquotesingle s area and the integral of the product of its two null expansions, using a fixed normalization. Trautman--Bondi mass is the total mass at a cut of future null infinity. A null Penrose inequality compares an appropriate initial horizon or section area with an asymptotic mass along a null hypersurface.

\paragraph{Target retained from the canonical source catalogue.}
\paragraph{Catalogue wording retained for traceability.}
For general null cones, characterize convergence of the Hawking energy, compare its limit with the Trautman–Bondi mass, and derive the corresponding null Penrose inequality.

\paragraph{Source scope and qualification.}
The cone foliation, expansion normalization, asymptotic hypotheses and choice of initial section must be kept from the report. A limit along one convenient foliation is not automatically a foliation-independent statement for all null cones.
\subsection{Short English statement}
Control Hawking-energy limits on general null cones, identify their relation to Bondi mass and derive the associated null Penrose bound.
\subsection{Sources}
\begin{itemize}
\item[\hypertarget{source:statement:1}{S1}] ULAM.AI and the UnsolvedMath Contributors, UnsolvedMath, catalogue entry. Catalogue curation under CC BY 4.0; definitions and English exposition added locally.
\url{https://huggingface.co/datasets/ulamai/UnsolvedMath}.
\item[\hypertarget{source:statement:2}{S2}] Mathematical General Relativity (2018), Oberwolfach Reports, DOI 10.4171/owr/2018/36, pp. 2157--2251. Primary report accessed 2026-10-11; relevant PDF pages 83, 84.
\url{https://doi.org/10.4171/owr/2018/36}.
\item[\hypertarget{source:statement:3}{S3}] Mars--Soria 2015 (catalogue-listed context; not independently consulted in this preparation).
\url{https://arxiv.org/abs/1506.01545}.
\item[\hypertarget{source:statement:4}{S4}] Mars--Soria null Penrose inequality (catalogue-listed context; not independently consulted in this preparation).
\url{https://arxiv.org/abs/1511.06242}.
\end{itemize}
\label{end:statement}
\end{document}
